\section{A sharp low-rank optimizer with a separable remainder}
\label{sec:low-rank-optimizer}

The preceding rank bound controls every point of the optimal face. A stronger
existence statement holds independently of the target's Schmidt rank. It also
controls the form of an attaining ensemble.

\begin{theorem}[A low-rank pure-target optimizer]
\label{thm:low-rank-pure-target}
Let $A,B$ be arbitrary finite-dimensional density matrices, let $x$ be any
vector in their bipartite space, and put
\[
                 D=\max\{\operatorname{rank}A,\operatorname{rank}B\}.
\]
Among the states with actual marginals $A,B$, the maximum of
$\langle x,\rho x\rangle$ is attained by a state of the form
\[
 \boxed{\quad
 \rho=ww^\dagger+\sum_{j=1}^{q}p_j
                 |a_j\otimes b_j\rangle\langle a_j\otimes b_j|,
 \qquad 0\le q\le D-1,\quad p_j\ge0,
 \quad}
\]
where the local vectors $a_j,b_j$ are unit vectors. Consequently
\[
                         \boxed{\operatorname{rank}\rho\le D.}
\]
The pure core is subnormalized: $\|w\|^2+\sum_jp_j=1$.
The theorem includes singular marginals, rectangular systems, and the case
in which the optimal target overlap is zero. If the remainder is nonzero,
its product vectors can be chosen with equal weights $p_j=\delta/q$.
\end{theorem}

We first give the elementary coupling construction used in the proof.

\begin{lemma}[A short separable coupling of arbitrary marginals]
\label{lem:fourier-separable-coupling}
Suppose $E,F\ge0$ have the same positive trace $\delta$, and let
$q=\max\{\operatorname{rank}E,\operatorname{rank}F\}$. There is a
positive separable operator $R$ with marginals $E,F$ that is a sum of $q$
product projectors, and therefore has rank at most $q$. No separable
coupling with these marginals can have fewer than $q$ nonzero product
terms.
\end{lemma}

\begin{proof}
Diagonalize the marginals and pad their positive eigenvalue lists with zeros
to length $q$:
\[
 E=\sum_{i=0}^{q-1}a_i|e_i\rangle\langle e_i|,
 \qquad F=\sum_{j=0}^{q-1}b_j|f_j\rangle\langle f_j|,
 \qquad \sum_i a_i=\sum_jb_j=\delta.
\]
Terms with zero coefficients are omitted if the corresponding physical
space has dimension less than $q$. Write $\omega=e^{2\pi i/q}$ and define
\[
 u_t=\sum_i\sqrt{a_i}\,\omega^{ti}e_i,
 \qquad v_t=\sum_j\sqrt{b_j}\,\omega^{-tj}f_j,
 \qquad t=0,\ldots,q-1.
\]
Both vectors have squared norm $\delta$. Finite Fourier orthogonality gives
\[
 {1\over q}\sum_tu_tu_t^\dagger=E,
 \qquad {1\over q}\sum_tv_tv_t^\dagger=F.
\]
Hence
\[
 R={1\over q\delta}\sum_{t=0}^{q-1}
       (u_tu_t^\dagger)\otimes(v_tv_t^\dagger)
   ={\delta\over q}\sum_{t=0}^{q-1}
       \left|{u_t\over\sqrt\delta}\otimes
                    {v_t\over\sqrt\delta}\right\rangle
       \left\langle{u_t\over\sqrt\delta}\otimes
                    {v_t\over\sqrt\delta}\right|
\]
has the required marginals and is explicitly separable with $q$ terms.
Conversely, each product projector contributes rank at most one to each
marginal, so any separable ensemble requires at least
$\max\{\operatorname{rank}E,\operatorname{rank}F\}$ terms.

For comparison with a cyclic-matching construction, define
\[
 z_\ell=\sum_{i=0}^{q-1}
       \sqrt{a_i b_{i+\ell}/\delta}\;e_i\otimes f_{i+\ell},
       \qquad \ell=0,\ldots,q-1,
\]
with indices modulo $q$. The same Fourier identity gives
$R=\sum_\ell z_\ell z_\ell^\dagger$: both sides have entry
$\sqrt{a_i b_j a_k b_l}/\delta$ precisely when
$j-i=l-k\pmod q$, and entry zero otherwise. Individual $z_\ell$ may be
entangled; their sum of projectors nevertheless has the displayed
separable decomposition. Its rank is exactly $q$, because the $q$
nonzero matching vectors have disjoint supports when one eigenvalue
list has $q$ positive entries.
\end{proof}

\begin{proof}[Proof of Theorem~\ref{thm:low-rank-pure-target}]
Every feasible state is supported on
$\operatorname{ran}A\otimes\operatorname{ran}B$. Compress to these
supports and compress $x$ accordingly. Write their dimensions as $m,n$;
then $D=\max(m,n)$. The compressed target need not be a unit vector and
may vanish. The proof below allows both possibilities.

Use the coefficient-matrix convention
$w=\sum_{ij}W_{ij}|i,j\rangle$, so the pure marginal matrices are
$WW^\dagger$ and $(W^\dagger W)^T$. Set $C=B^T$ and consider
\[
 \mathcal K=\{W:\ WW^\dagger\le A,\quad W^\dagger W\le C\}.
\]
This is a nonempty compact set: it contains zero, is closed, and every
element obeys $\|W\|_F^2\le1$. It is convex, for example by the two
equivalent Schur-complement linear matrix inequalities. If $X$ is the
coefficient matrix of the compressed target, define
\[
 \tau=\max_{W\in\mathcal K}\operatorname{Re}\operatorname{Tr}(X^\dagger W).
\]
Phase invariance of $\mathcal K$ shows that this is also
$\max_{W\in\mathcal K}|\operatorname{Tr}(X^\dagger W)|$, and the exact
pure-target reduction identifies the optimal state overlap as $\tau^2$.
For completeness, the two directions of that reduction are the
rank-one extraction $w=\rho x/\sqrt{\langle x,\rho x\rangle}$ when the
denominator is nonzero, and the product-residual completion of any
feasible $w$. A zero state overlap uses $w=0$ instead. These arguments
do not require an invertible target or an invertible coupling.

Among all maximizers of the real linear objective, choose $W$ maximizing
$\|W\|_F^2$. Such a choice exists because the maximizing set is compact.
Put
\[
 E=A-WW^\dagger,\qquad F=C-W^\dagger W,
 \qquad \delta=\operatorname{Tr}E=\operatorname{Tr}F=1-\|W\|_F^2.
\]
If $\delta=0$, both residual matrices vanish, so $ww^\dagger$ is already
an optimal state of rank one. Suppose henceforth that $\delta>0$.

We claim
\[
                 \max\{\operatorname{rank}E,\operatorname{rank}F\}
                    \le D-1.                                      \tag{LR1}
\]
First assume $m\ge n$. If $E$ were positive definite, let
$V=\operatorname{ran}F$, of dimension $s\ge1$. Since
\[
                 \dim W(V^\perp)\le n-s<m,
\]
there is a nonzero vector $u\perp W(V^\perp)$. Choose any nonzero
$v\in V$ and set $Y=uv^\dagger$. Then $W^\dagger u\in V$, so all
three matrices
$W^\dagger Y$, $Y^\dagger W$, and $Y^\dagger Y$ are supported on $V$.
For sufficiently small $\varepsilon>0$, both signs satisfy
\[
 \begin{split}
 A-(W\pm\varepsilon Y)(W\pm\varepsilon Y)^\dagger
    &=E\mp\varepsilon(WY^\dagger+YW^\dagger)
                           -\varepsilon^2YY^\dagger\ge0,\\
 C-(W\pm\varepsilon Y)^\dagger(W\pm\varepsilon Y)
    &=F\mp\varepsilon(W^\dagger Y+Y^\dagger W)
                           -\varepsilon^2Y^\dagger Y\ge0.
 \end{split}
\]
The first assertion follows from strict positivity of $E$ on its full
space. The second follows from strict positivity of $F$ on $V$, since
the perturbations vanish on $V^\perp$ and have no cross blocks. Thus
both $W+\varepsilon Y$ and $W-\varepsilon Y$ lie in $\mathcal K$.
Their linear objectives average to the optimal value, so both are
themselves maximizers. But
\[
 {\|W+\varepsilon Y\|_F^2+\|W-\varepsilon Y\|_F^2\over2}
       =\|W\|_F^2+\varepsilon^2\|Y\|_F^2>\|W\|_F^2,
\]
contradicting the choice of $W$. Therefore
$\operatorname{rank}E\le m-1$. If $m>n$, also
$\operatorname{rank}F\le n\le m-1$. If $m=n$, apply the same argument
to $W^\dagger$ with the two caps exchanged to obtain
$\operatorname{rank}F\le n-1$. The case $n>m$ is symmetric. This proves
(LR1).

The product completion
\[
                ww^\dagger+{E\otimes F^T\over\delta}
\]
has marginals $A,B$. Its pure term already has target overlap $\tau^2$,
the global optimum. Consequently
\[
       \langle x,E\otimes F^T\,x\rangle=0,
       \qquad
       x\perp\operatorname{ran}E\otimes\operatorname{ran}F^T.       \tag{LR2}
\]
Here the second conclusion follows because the positive product matrix
is strictly positive on its support. The argument remains valid for
$\tau=0$: the product completion cannot have positive target overlap
when zero is the global maximum.

Apply Lemma~\ref{lem:fourier-separable-coupling} to $E,F^T$. It gives a
separable residual $R$ with these actual marginals and at most
$q=\max(\operatorname{rank}E,\operatorname{rank}F)\le D-1$ product
terms. It is supported on the space in (LR2), hence
$\langle x,Rx\rangle=0$. Therefore $\rho=ww^\dagger+R$ has the correct
actual marginals and still attains the global target optimum. Its rank
is at most $1+q\le D$, with the asserted separable decomposition.
\end{proof}

\paragraph{Sharpness in every rectangular dimension.}
The bound $D$ cannot be lowered, even for positive definite marginals.
Suppose $m\ge n$, let $A=I_m/m$, and choose a positive definite $B$ with
an eigenvalue $\beta>(m-1)/m$ and unit eigenvector $b$. For any coupling,
the positive principal block
$T=(I\otimes\langle b|)\rho(I\otimes|b\rangle)$ satisfies
$0\le T\le A$ and $\operatorname{Tr}T=\beta$. Consequently
\[
       \beta\le {\operatorname{rank}T\over m}
             \le {\operatorname{rank}\rho\over m},
       \qquad \operatorname{rank}\rho\ge m.
\]
Such $B$ exists for every $n$: for $n>1$, take the remaining eigenvalues
equal to $\varepsilon$ with
$0<\varepsilon<1/[m(n-1)]$; for $n=1$, take $B=[1]$. The case $n>m$
is symmetric. This example forces rank $D$ for every coupling, and
therefore makes the theorem sharp for every choice of target.

An explicit unique optimizer of rank $m$ is also available for a
maximal-Schmidt-rank target in every $m\ge n$ with $m>1$. Let
$0<\delta<1$, $\mu=(1-\delta)/n$, and
\[
 x={1\over\sqrt n}\sum_{i=0}^{n-1}|i,i\rangle,
 \quad w=\sqrt\mu\sum_{i=0}^{n-1}|i,i\rangle,
 \quad E={\delta\over m-1}\sum_{i=1}^{m-1}|i\rangle\langle i|.
\]
Take
\[
 A=\mu\sum_{i=0}^{n-1}|i\rangle\langle i|+E,
 \qquad B=\mu I_n+\delta|0\rangle\langle0|,
 \qquad \rho=ww^\dagger+E\otimes|0\rangle\langle0|.
\]
Both actual marginals are positive definite. Each target-diagonal
population is capped by $\mu$, so positivity and Cauchy--Schwarz give
target overlap at most $n\mu=1-\delta$, attained by this $\rho$.
For $n>1$, equality forces every target population to equal $\mu$ and
the target-diagonal principal block to be $ww^\dagger$. All remaining
nonzero populations are confined to
$\operatorname{span}\{|i,0\rangle:i\ne0\}$. Positivity therefore
writes any optimizer as
$ww^\dagger+wy^\dagger+yw^\dagger+R$ on these two blocks.
The off-diagonal Alice marginal entries with index zero force $y=0$;
the Alice marginal then forces $R=E\otimes|0\rangle\langle0|$.
Thus the displayed state is the unique optimizer and has rank $m$.
For $n=1$, the only coupling is $A\otimes B$, also of rank $m$.
For $m=n=1$, the only state has rank one.

\begin{corollary}[Fully entangled fraction and other target orbits]
For fixed actual marginals $A,B$ on $\mathbb C^d\otimes\mathbb C^d$,
the largest possible fully entangled fraction is attained by a state of
rank at most $D=\max(\operatorname{rank}A,\operatorname{rank}B)$, with
one pure core and at most $D-1$ product residual terms. The same holds
when the maximally entangled targets are replaced by any compact family
of pure targets, including a local-unitary orbit with prescribed Schmidt
coefficients.
\end{corollary}

\begin{proof}
The coupling set and the target family are compact, so choose a joint
maximizer $(\rho_*,x_*)$. Apply the theorem to the fixed target $x_*$.
The resulting low-rank state attains the joint optimal value at that
target, and therefore attains the maximum over the target family as
well. No assumption about uniqueness of the optimizing target is used.
\end{proof}

\begin{corollary}[A shorter simultaneous coherent-minimum witness]
Let $\alpha,\beta$ be probability vectors of length $d$, set
$m_i=\min(\alpha_i,\beta_i)$, $u=\alpha-m$, $v=\beta-m$, and
$\delta=\sum_i u_i=\sum_i v_i$. There is one state with actual diagonal
marginals $\operatorname{diag}\alpha,\operatorname{diag}\beta$ that
simultaneously attains
\[
 \langle x_p,\rho x_p\rangle
       =\left(\sum_i\sqrt{p_i m_i}\right)^2,
 \qquad x_p=\sum_i\sqrt{p_i}|i,i\rangle,
\]
for every probability vector $p$. If $\delta=0$, this state is pure.
Otherwise it may be chosen with rank at most
\[
           1+\max\{|\{i:u_i>0\}|,|\{i:v_i>0\}|\}\le d
\]
and with a separable residual having this same maximum number of product
terms. The displayed rank is exact when $\sum_i m_i>0$; when all
$m_i=0$, the pure core vanishes and the rank is one smaller.
\end{corollary}

\begin{proof}
Put $z=\sum_i\sqrt{m_i}|i,i\rangle$. If $\delta>0$, replace the
product residual of the coherent-minimum construction by the Fourier
coupling $R$ of $\operatorname{diag}u,\operatorname{diag}v$ from
Lemma~\ref{lem:fourier-separable-coupling}. Their positive supports are
disjoint, so $R$ is supported off the target diagonal and annihilates
every $x_p$. Thus $\rho=zz^\dagger+R$ has the displayed overlaps and
the required actual marginals. If $z\ne0$, its range is orthogonal to
the range of $R$, so the ranks add exactly. If $z=0$, all the overlaps
vanish and the same construction has rank equal to the residual rank,
which is an even stronger bound. Both residual supports are nonempty
when $\delta>0$ and are disjoint subsets of $\{1,\ldots,d\}$, so their
maximum size is at most $d-1$. If $\delta=0$, use $\rho=zz^\dagger$.
The exact target-population inequality proves optimality for each $p$;
when the spectra and weights have the required order, this also gives
a shorter attaining witness for the spectral-overlap theorem.
\end{proof}

\paragraph{Prior-art boundary and verification scope.}
The existence of some coupling with rank at most the larger marginal
rank is already Theorem~6.4.1 of A.~Klyachko,
\emph{Quantum marginal problem and representations of the symmetric
group}, arXiv:quant-ph/0409113. That theorem also discusses simultaneous
maximization of the spectral norm. The elementary Fourier ensemble
construction above is not asserted to be new. The additional statement
here is simultaneous maximization of an arbitrary prescribed pure-target
overlap, with a pure core and at most $D-1$ product residual terms.
Targeted searches of primary quantum-coupling, marginal, and quantum
transport sources did not locate this exact statement; that bounded
search does not establish publication priority. The perturbation and
rank arguments in this section are exact mathematical proofs, separate
from the Lean-verified extraction/completion reduction. No Lean
formalization of these new rank arguments is claimed.
